% TOP_PROBLEM: {"id": 30006925, "title": "Absolute continuity of Bernoulli convolutions outside reciprocal Pisot parameters", "category_id": 19, "category": {"id": 19, "name": "probability", "display_name": "Probability", "description": "Problems involving probability theory, stochastic processes, and random structures.", "slug": "probability", "order_index": 19, "created_at": "2026-07-31T15:26:25.671Z"}, "status": "open"}
% FIRST_POSTED: 2026-09-27
% !TeX program = lualatex
% ATTEMPT_STATUS: unresolved
\documentclass[11pt]{article}
\usepackage[margin=25mm]{geometry}
\usepackage{fontspec}
\setmainfont{Latin Modern Roman}
\usepackage{amsmath,amssymb,mathtools}
\usepackage{unicode-math}
\setmathfont{Latin Modern Math}
\usepackage{xurl,hyperref}
\hypersetup{colorlinks=true,urlcolor=blue}
\setcounter{secnumdepth}{0}
\setlength{\parindent}{0pt}
\setlength{\parskip}{5pt}
\setlength{\emergencystretch}{3em}
\begin{document}
\section*{\texorpdfstring{Absolute continuity of Bernoulli convolutions outside reciprocal Pisot parameters}{Absolute continuity of Bernoulli convolutions outside reciprocal Pisot parameters}}\label{30006925}
\subsection{Definitions and mathematical statement}
For $1/2<\lambda<1$, let $(\epsilon_n)_{n\ge0}$ be independent random signs, each equally likely to be $1$ or $-1$, and let $\nu_\lambda$ be the probability distribution of the almost surely convergent series
\[
 X_\lambda=\sum_{n=0}^{\infty}\epsilon_n\lambda^n.
\]
A Pisot number is a real algebraic integer $\beta>1$ all of whose other algebraic conjugates have modulus less than one. The selected assertion is
\[
 \lambda^{-1}\text{ is not Pisot}\quad\Longrightarrow\quad
 \nu_\lambda\ll\mathcal L^1.
\]
Absolute continuity means that every Lebesgue-null Borel set has $\nu_\lambda$-measure zero, equivalently that $\nu_\lambda$ has an integrable density. The quantifier is over every eligible parameter, not almost every parameter. Full Hausdorff dimension alone would not establish absolute continuity. Multiplying $X_\lambda$ by a nonzero normalization constant leaves this question unchanged.
\par\smallskip\textit{Formulation and concept references:} \href{https://londmathsoc.onlinelibrary.wiley.com/doi/full/10.1112/jlms.70515}{[S1]}.

\subsection{Short English statement}
Must a random signed geometric series have a probability density whenever the reciprocal of its contraction parameter is not a Pisot number?
\subsection{Research attempt}
\paragraph{Process and principal gap.}
This attempt was generated by GPT-6 Astra Ultra on 27 September 2026. We prove absolute continuity for an explicit infinite family of non-Pisot reciprocal parameters, compute their densities, and test Fourier decay and full dimension as proposed routes to the general statement. The first unresolved step is obtaining sufficient quantitative regularity for an arbitrary fixed parameter whose reciprocal is not Pisot.

The introduction of [S1] explicitly distinguishes the deterministic Bernoulli convolution from its random-parameter model. It records almost-everywhere absolute continuity, small exceptional parameter sets, and full-dimension results without concluding the pointwise assertion here. Those qualifications matter: a small exceptional set can still contain a specified parameter. We use [S1] for this literature check, while the constructions and estimates below are proved directly.

\paragraph{A family that splits into independent uniform variables.}
Fix an integer $m\ge2$ and put $\lambda=2^{-1/m}$. Splitting the defining series into residue classes modulo $m$ gives an almost surely absolutely convergent identity
\[
 X_\lambda=\sum_{r=0}^{m-1}\lambda^rY_r,
 \qquad Y_r=\sum_{j=0}^{\infty}\varepsilon_{mj+r}2^{-j}.
 \tag{426.1}
\]
The variables $Y_r$ are independent because their sign sequences are disjoint. If $\varepsilon_j=2b_j-1$ with independent fair bits $b_j$, then
\[
 \sum_{j\ge0}\varepsilon_j2^{-j}
 =4\sum_{j\ge0}b_j2^{-(j+1)}-2.
\]
The binary series on the right is uniform on $[0,1]$: every dyadic interval of length $2^{-n}$ has probability $2^{-n}$, and the countable set of ambiguous binary endpoints has probability zero. Hence each $Y_r$ is uniform on $[-2,2]$.

It follows that $\nu_\lambda$ has density
\[
 p_\lambda=f_0*f_1*\cdots*f_{m-1},\qquad
 f_r(x)=\frac{1}{4\lambda^r}1_{[-2\lambda^r,2\lambda^r]}(x).
 \tag{426.2}
\]
Every $f_r$ has integral one. Convolution with a probability density does not increase an $L^\infty$ norm, so $\|p_\lambda\|_\infty\le\|f_0\|_\infty=1/4$. The support is
\[
 \left[-2\sum_{r=0}^{m-1}\lambda^r,\,
        2\sum_{r=0}^{m-1}\lambda^r\right]
 =[-(1-\lambda)^{-1},(1-\lambda)^{-1}],
\]
in agreement with the maximum possible absolute value of the original series. No limiting density approximation is required for this family.

\paragraph{An explicit formula and verification of the parameter restriction.}
Set $\ell_r=4\lambda^r$ and $S=2\sum_r\lambda^r$. Translating the uniform variables to $[0,\ell_r]$ and computing the volume of a clipped simplex gives
\[
 p_\lambda(x)=\frac{1}{(m-1)!\prod_r\ell_r}
 \sum_{J\subseteq\{0,\ldots,m-1\}}(-1)^{|J|}
       \left(x+S-\sum_{r\in J}\ell_r\right)_+^{m-1}.
 \tag{426.3}
\]
Indeed the distribution function is the volume of $\{u_r\ge0:\sum_ru_r\le x+S,\ u_r\le\ell_r\}$ divided by $\prod_r\ell_r$. Inclusion-exclusion removes the upper-coordinate violations; the unrestricted simplex has volume $(x+S)_+^m/m!$. Differentiating gives (426.3). Thus the density is a continuous piecewise polynomial, and in fact belongs to $C^{m-2}$. One can also see the regularity by differentiating $m-2$ distinct interval factors distributionally: the remaining convolution of two interval densities is continuous.

The reciprocal $\beta=2^{1/m}$ is an algebraic integer with irreducible polynomial $X^m-2$, by Eisenstein's criterion at $2$. All its conjugates have modulus $\beta>1$. Since $m\ge2$, at least one other conjugate exists and lies outside the unit disk. Therefore $\beta$ is not Pisot. Also $1/2<\lambda<1$. Equations (426.1)--(426.3) genuinely apply to parameters in the target class, rather than only to excluded or endpoint cases.

\paragraph{The Fourier product and a checked excluded parameter.}
With convention $\widehat\nu(t)=\int e^{itx}\,d\nu(x)$, independence gives
\[
 \widehat\nu_\lambda(t)=\prod_{n=0}^{\infty}\cos(\lambda^nt).
 \tag{426.4}
\]
For each fixed $t$, the product converges because $\sum_n|1-\cos(\lambda^nt)|<\infty$. Absolute continuity would imply $\widehat\nu_\lambda(t)\to0$ as $|t|\to\infty$ by the Riemann--Lebesgue lemma.

Consider the excluded golden-ratio case $\beta=(1+\sqrt5)/2$, $\lambda=\beta^{-1}$. Its other conjugate $\alpha=(1-\sqrt5)/2$ satisfies $|\alpha|<1$, and $\beta^j+\alpha^j$ is an integer for every nonnegative integer $j$. Thus $\beta^j$ approaches an integer exponentially as $j\to+\infty$, while $\beta^j\to0$ as $j\to-\infty$. It follows that
\[
 \sum_{j\in\mathbb Z}\bigl(1-|\cos(2\pi\beta^j)|\bigr)<\infty.
\]
None of these cosine factors is zero. For $j\ne0$, $\beta^j$ is irrational: a rational value would equal its conjugate $\alpha^j$, contradicting their different absolute values. For $j=0$, the cosine is one. An infinite product of these absolute values is therefore strictly positive, say $c>0$. At $t_k=2\pi\beta^k$, (426.4) gives
\[
 |\widehat\nu_\lambda(t_k)|
 =\prod_{j\le k}|\cos(2\pi\beta^j)|\ge c. \tag{426.5}
\]
So this particular measure is not absolutely continuous. The argument verifies why the Pisot restriction is meaningful; it does not give a counterexample within the permitted non-Pisot class.

For a general parameter, the failure of this specific concentration argument is not a proof of absolute continuity. Negating one sufficient obstruction does not establish a sufficient regularity estimate. Even a proof of decay to zero would be weaker than the integrability of the Fourier transform that a direct density argument could use.

\paragraph{A sufficient quantitative route, with its assumption exposed.}
If one could prove
\[
 \int_{\mathbb R}\prod_{n=0}^{\infty}|\cos(\lambda^nt)|^2\,dt<\infty,
 \tag{426.6}
\]
then $\nu_\lambda$ would have an $L^2$ density. For completeness, convolve the measure with a Gaussian approximate identity $\phi_\varepsilon$. Plancherel and $|\widehat\phi_\varepsilon|\le1$ give a uniform $L^2$ bound on $\nu_\lambda*\phi_\varepsilon$ from (426.6). A weakly convergent subsequence in $L^2$ has a density limit; testing against smooth compactly supported functions identifies that limit with the original measure. This proves the asserted sufficient criterion.

No estimate in this attempt establishes (426.6) for all the desired parameters. It is also stronger than the original request, which asks only for absolute continuity, not an $L^2$ density. Thus failure to prove this particular criterion would not refute the conjecture.

\paragraph{Full dimension is not the missing density estimate.}
Here is a concrete measure illustrating the logical distinction without making any claim that it is a Bernoulli convolution. On $[0,1]$, choose independent fair binary digits except that the digits in positions $1,4,9,16,\ldots$ are fixed to zero. Let $\mu$ be the resulting law. The set of numbers satisfying the first $r$ fixed-digit restrictions has Lebesgue measure $2^{-r}$, so their intersection is a null set supporting $\mu$.

At binary depth $n$, an allowed cylinder has mass $2^{-(n-\lfloor\sqrt n\rfloor)}$. For every $s<1$ this is at most $C_s2^{-sn}$. An interval whose length lies between $2^{-(n+1)}$ and $2^{-n}$ intersects at most two depth-$n$ cylinders, apart from endpoints of zero mass. Hence $\mu(I)\le C'_s|I|^s$. Covering any set of full $\mu$-measure by intervals now shows its Hausdorff dimension is at least $s$; letting $s\uparrow1$ gives dimension one. Nevertheless $\mu$ is singular by the preceding null-support calculation.

Thus a full-dimension theorem, even for a measure with no atoms, does not generally imply absolute continuity. Further structure of Bernoulli convolutions would have to supply that additional implication; it is not contained in the numerical dimension value alone.

\paragraph{Verification and outcome.}
Finite exact checks verify the binary grouping and the inclusion-exclusion density formula in small cases; the infinite product and full-dimension arguments are justified analytically above. Source statements were checked for the differences between deterministic parameters, random parameter sequences, almost-everywhere assertions, and pointwise assertions.

\textbf{UNRESOLVED.} The family $\lambda=2^{-1/m}$, $m\ge2$, has explicitly computed bounded densities and non-Pisot reciprocals. The general target remains open in this attempt because no sufficient density or Fourier-integrability estimate is established for every remaining allowed parameter.

\subsection{Sources}
\begin{itemize}
\item[S1]Baker, Koivusalo, Troscheit and Zhang, On the Fourier transform of random Bernoulli convolutions, JLMS 113 (2026), e70515.
\url{https://londmathsoc.onlinelibrary.wiley.com/doi/full/10.1112/jlms.70515}.
\textit{Formulation source.}
\end{itemize}
\end{document}
